\section{How the work was done}\label{app:how}

\emph{Disclosure of generative-AI use.} This paper and its Lean development were produced substantially by a generative AI system, Claude
(Anthropic), working under the direction of the named author, who takes full responsibility for its contents. No AI system is listed as an author.

The local inequalities of Section~\ref{sec:cert} were found by linear programming over gap vectors located by floating-point search and
verified by interval branch-and-bound; the certificates were re-run, some box for box, and their inputs and closing constants were checked
again with independently written code. What is checked by the Lean kernel, and under which hypotheses, is stated in \S\ref{ssec:formal}. The
Lean libraries, the certificate data with their certifiers and logs, the generator of the replay files, and the build logs and
\texttt{\#print axioms} outputs of the release are in the repository \repourl{} at tag \repotag{} (Appendix~\ref{app:layout}).

\emph{Scope of the verification.} Theorems~\ref{thm:main} and~\ref{thm:S} and Corollary~\ref{cor:SC} are Lean theorems that depend only on
the three standard axioms. For the $K=5$ pair of Theorem~\ref{thm:main} the local certificate is replayed in the Lean kernel, so these two
statements have no hypothesis. For the $K=7$ pair, Theorem~\ref{thm:S} and Corollary~\ref{cor:SC} the local certificate is a displayed
hypothesis: a finite inequality checked outside Lean by interval arithmetic. The paper has not been refereed, and the Lean development has
not been rebuilt outside this project.
